% ============================================================
% Bảng liệt kê 100 câu hỏi của bộ đánh giá CheoBench v2
%
% BẮT BUỘC thêm vào preamble của main document:
%     \usepackage{longtable}   % cho bảng đa trang
%
% Cross-reference: \ref{tab:cheobench_full}
% ============================================================

\chapter*{Phụ lục}
\begin{longtable}{|p{1.9cm}|p{5.2cm}|p{7.3cm}|}
  \hline
  \textbf{ID} & \textbf{Câu hỏi} & \textbf{Đáp án kỳ vọng} \\
  \hline
  \endfirsthead
  \hline
  \textbf{ID} & \textbf{Câu hỏi} & \textbf{Đáp án kỳ vọng} \\
  \hline
  \endhead
  \hline
  \endfoot
  \hline
  \endlastfoot

  \texttt{CASE\_001} & Mô tả đặc điểm nhân vật Thị Kính? & Con gái Mãng Ông, bị oan khuất phải lên chùa tu hành dưới danh Kính Tâm, chịu đựng mọi nỗi oan trái với lòng ngay thẳng. \\
  \hline
  \texttt{CASE\_002} & Mô tả đặc điểm nhân vật Thị Màu? & Con gái phú ông, táo bạo, hồn nhiên; cưa cẩm chú tiểu Kính Tâm không được, đổ oan thai nhi cho Thị Kính. \\
  \hline
  \texttt{CASE\_003} & Mô tả đặc điểm nhân vật Súy Vân? & Con gái huyện Tể, gả cho Kim Nham. Bị Trần Phương xúi giục, giả điên để thoát khỏi chồng. \\
  \hline
  \texttt{CASE\_004} & Mô tả đặc điểm nhân vật Châu Long? & Vợ ba của Dương Lễ. Âm thầm nuôi dưỡng Lưu Bình ăn học suốt ba năm, giữ trọn tiết hạnh. \\
  \hline
  \texttt{CASE\_005} & Mô tả đặc điểm nhân vật Lưu Bình? & Người tài hoa nhưng lận đận. Nhờ Châu Long dìu dắt, đỗ Trạng nguyên. \\
  \hline
  \texttt{CASE\_006} & Mô tả đặc điểm nhân vật Trương Viên? & Người chồng thủy chung, con hiếu nghĩa. Sau 18 năm xa quê, nhờ tiếng hát xẩm nhận ra vợ. \\
  \hline
  \texttt{CASE\_007} & Mô tả đặc điểm nhân vật Thị Phương? & Vợ Trương Viên, khoét mắt mình cho mẹ chồng ăn. Gặp lại chồng nhờ đôi ngọc lưu ly. \\
  \hline
  \texttt{CASE\_008} & Mô tả đặc điểm nhân vật Đào Huế? & Vợ cũ Chu Mãi Thần, bỏ chồng theo Tuần Ty khi chồng còn nghèo khổ. \\
  \hline
  \texttt{CASE\_009} & Mô tả đặc điểm nhân vật Thiệt Thê? & Nhân vật nông nổi, ham vật chất, phải trả giá cho sai lầm của mình. \\
  \hline
  \texttt{CASE\_010} & Mô tả đặc điểm nhân vật Tuần Ty? & Viên quan tằng tịu với vợ cũ Chu Mãi Thần, định lấy làm vợ hai. \\
  \hline
  \texttt{CASE\_011} & Ai đóng vai Thị Kính trong vở Quan Âm Thị Kính? & Thảo Hiền và Hồng Thắm \\
  \hline
  \texttt{CASE\_012} & Ai đóng vai Thị Màu trong vở Quan Âm Thị Kính? & Vân Quyền và An Chinh \\
  \hline
  \texttt{CASE\_013} & Ai đóng vai Súy Vân trong vở Kim Nham? & Tạ Thị Kim Liên và Thúy Ngần \\
  \hline
  \texttt{CASE\_014} & Ai đóng vai Sùng Bà trong vở Quan Âm Thị Kính? & Thanh Ngoan và Trần Thị Thân \\
  \hline
  \texttt{CASE\_015} & Ai đóng vai Châu Long trong vở Lưu Bình Dương Lễ? & Ngọc Ánh và Phương Mây \\
  \hline
  \texttt{CASE\_016} & Ai đóng vai Đào Huế trong vở Chu Mãi Thần? & Thanh Ngoan và Kim Liên \\
  \hline
  \texttt{CASE\_017} & Ai đóng vai Từ Thức trong vở Từ Thức? & Phú Kiên \\
  \hline
  \texttt{CASE\_018} & Ai đóng vai Lưu Bình trong vở Lưu Bình Dương Lễ? & Nguyễn Duy và Hồng Nam \\
  \hline
  \texttt{CASE\_019} & Vân Quyền đóng vai gì trong vở Quan Âm Thị Kính? & Vân Quyền đóng vai Thị Màu \\
  \hline
  \texttt{CASE\_020} & Huy Toàn đóng vai gì trong vở Quan Âm Thị Kính? & Huy Toàn đóng vai Thiện Sỹ \\
  \hline
  \texttt{CASE\_021} & Tuấn Kha đóng vai gì trong vở Quan Âm Thị Kính? & Tuấn Kha đóng vai Mãng Ông (bố Thị Kính) \\
  \hline
  \texttt{CASE\_022} & Bá Dũng đóng vai gì trong vở Kim Nham? & Bá Dũng đóng vai Trần Phương \\
  \hline
  \texttt{CASE\_023} & Trần Hải đóng vai gì trong vở Kim Nham? & Trần Hải đóng vai Phù thủy trong trích đoạn Phù thủy sợ ma \\
  \hline
  \texttt{CASE\_024} & An Chinh đóng vai gì trong vở Kim Nham? & An Chinh đóng vai Hỷ đồng \\
  \hline
  \texttt{CASE\_025} & Ngọc Minh đóng vai gì trong vở Chu Mãi Thần? & Ngọc Minh đóng vai Tuần Ty \\
  \hline
  \texttt{CASE\_026} & Vở Quan Âm Thị Kính có những trích đoạn nào? & Cắt râu, Lý trưởng - Mẹ Mõ, Mầu - Nô - Phú Ông, Thị Mầu lên chùa, Vu quy \\
  \hline
  \texttt{CASE\_027} & Trích đoạn 'Thị Mầu lên chùa' thuộc vở nào? & Trích đoạn Thị Mầu lên chùa thuộc vở Quan Âm Thị Kính \\
  \hline
  \texttt{CASE\_028} & Nhân vật chính trong trích đoạn 'Thị Mầu lên chùa' là ai? & Nhân vật chính là Thị Màu và Thị Kính (Kính Tâm) \\
  \hline
  \texttt{CASE\_029} & Trích đoạn 'Vu quy' trong vở Quan Âm Thị Kính miêu tả sự kiện gì? & Vu quy là cảnh đám cưới của Thị Kính và Thiện Sỹ \\
  \hline
  \texttt{CASE\_030} & Vở Kim Nham có những trích đoạn nào? & Kim Nham có 4 trích đoạn: Mụ quán - Trần Phương, Phù thủy sợ ma, Súy Vân giả dại, Trần Phương vào chùa \\
  \hline
  \texttt{CASE\_031} & Trích đoạn 'Súy Vân giả dại' thuộc vở nào? & Thuộc vở Kim Nham \\
  \hline
  \texttt{CASE\_032} & Trích đoạn 'Lớp Tiên Nữ - Đoàn tụ' thuộc vở nào? & Thuộc vở Trương Viên \\
  \hline
  \texttt{CASE\_033} & Khắc Huy đóng vai gì trong vở Trương Viên? & Khắc Huy đóng vai Hề trong trích đoạn Lớp Tiên Nữ - Đoàn tụ \\
  \hline
  \texttt{CASE\_034} & Ngọc Ánh đóng vai gì trong vở Lưu Bình Dương Lễ? & Ngọc Ánh đóng vai Châu Long \\
  \hline
  \texttt{CASE\_035} & Mạnh Phóng đóng vai gì trong vở Quan Âm Thị Kính? & Mạnh Phóng đóng vai Nô \\
  \hline
  \texttt{CASE\_036} & Tất cả những diễn viên nào đã đóng vai Thị Màu? & Vân Quyền và An Chinh \\
  \hline
  \texttt{CASE\_037} & Tất cả những diễn viên nào đã đóng vai Thị Kính? & Thảo Hiền và Hồng Thắm \\
  \hline
  \texttt{CASE\_038} & Tất cả những diễn viên nào đã đóng vai Súy Vân? & Tạ Thị Kim Liên và Thúy Ngần \\
  \hline
  \texttt{CASE\_039} & Tất cả những diễn viên nào đã đóng vai Sùng Bà? & Thanh Ngoan và Trần Thị Thân \\
  \hline
  \texttt{CASE\_040} & Tất cả những diễn viên nào đã đóng vai Thiện Sỹ? & Huy Toàn và Trần Xuân Tài \\
  \hline
  \texttt{CASE\_041} & Tất cả những diễn viên nào đã đóng vai Hỷ đồng? & An Chinh và Thanh Mai \\
  \hline
  \texttt{CASE\_042} & Tất cả những diễn viên nào đã đóng vai Châu Long? & Ngọc Ánh và Phương Mây \\
  \hline
  \texttt{CASE\_043} & Tất cả những diễn viên nào đã đóng vai Tuần Ty? & Ngọc Minh và Phú Kiên \\
  \hline
  \texttt{CASE\_044} & Tất cả những diễn viên nào đã đóng vai Đào Huế? & Thanh Ngoan và Kim Liên \\
  \hline
  \texttt{CASE\_045} & Liệt kê tất cả nhân vật trong vở Quan Âm Thị Kính? & Thị Kính, Thị Màu, Thiện Sỹ, Sùng Bà, Sùng Ông, Nô, Mãng Ông, Phú Ông, Lý Trưởng, Mẹ Mõ \\
  \hline
  \texttt{CASE\_046} & Liệt kê tất cả nhân vật trong vở Kim Nham? & Súy Vân, Trần Phương, Mụ Quán, Hỷ đồng, Khoèo, Phù thủy \\
  \hline
  \texttt{CASE\_047} & Liệt kê tất cả nhân vật trong vở Chu Mãi Thần? & Đào Huế, Thiệt Thê, Tuần Ty \\
  \hline
  \texttt{CASE\_048} & Liệt kê tất cả nhân vật trong vở Lưu Bình Dương Lễ? & Lưu Bình, Châu Long, Dương Lễ \\
  \hline
  \texttt{CASE\_049} & Những diễn viên nào đã tham gia vở Quan Âm Thị Kính? & Thảo Hiền, Hồng Thắm, Vân Quyền, An Chinh, Huy Toàn, Trần Xuân Tài, Thanh Ngoan, Trần Thị Thân, Bá Dũng, Ngọc Minh, Văn Quân, Mạnh Phóng, Tuấn Kha \\
  \hline
  \texttt{CASE\_050} & Những diễn viên nào đã tham gia vở Kim Nham? & Tạ Thị Kim Liên, Thúy Ngần, An Chinh, Thanh Mai, Đăng Toàn, Tuấn Kha, Trần Hải, Bá Dũng, Huy Toàn, Thanh Mạn, Ngọc Minh \\
  \hline
  \texttt{CASE\_051} & Những diễn viên nào đã tham gia đồng thời cả vở Quan Âm Thị Kính và Kim Nham? & An Chinh, Tuấn Kha, Bá Dũng, Huy Toàn, Ngọc Minh \\
  \hline
  \texttt{CASE\_052} & Mối quan hệ giữa Thị Kính và Thiện Sỹ là gì? & Thị Kính là vợ của Thiện Sỹ trong vở Quan Âm Thị Kính \\
  \hline
  \texttt{CASE\_053} & Mối quan hệ giữa Sùng Bà và Thiện Sỹ là gì? & Sùng Bà là mẹ của Thiện Sỹ \\
  \hline
  \texttt{CASE\_054} & Mối quan hệ giữa Châu Long và Dương Lễ là gì? & Châu Long là vợ ba của Dương Lễ \\
  \hline
  \texttt{CASE\_055} & Mối quan hệ giữa Lưu Bình và Dương Lễ là gì? & Lưu Bình và Dương Lễ là bạn thân từ thời đi học \\
  \hline
  \texttt{CASE\_056} & Mối quan hệ giữa Trương Viên và Thị Phương là gì? & Trương Viên là chồng của Thị Phương \\
  \hline
  \texttt{CASE\_057} & Trần Phương tác động đến Súy Vân như thế nào? & Trần Phương xúi giục Súy Vân giả điên để bỏ chồng Kim Nham \\
  \hline
  \texttt{CASE\_058} & Thị Phương hy sinh điều gì vì Trương Mẫu? & Thị Phương khoét mắt mình để cho mẹ chồng Trương Mẫu ăn khi đói khát \\
  \hline
  \texttt{CASE\_059} & Những nhân vật nào trong Quan Âm Thị Kính tương tác trực tiếp với Thị Màu? & Thị Kính (Kính Tâm), Nô, Phú Ông, Lý Trưởng, Mẹ Mõ \\
  \hline
  \texttt{CASE\_060} & Thị Kính bị Sùng Bà vu oan như thế nào? & Sùng Bà vu oan Thị Kính định giết chồng khi cắt râu cho Thiện Sỹ, đuổi nàng khỏi nhà \\
  \hline
  \texttt{CASE\_061} & Có bao nhiêu diễn viên khác nhau từng đóng vai Thị Màu? & 2 diễn viên: Vân Quyền và An Chinh \\
  \hline
  \texttt{CASE\_062} & Có bao nhiêu diễn viên khác nhau từng đóng vai Thị Kính? & 2 diễn viên: Thảo Hiền và Hồng Thắm \\
  \hline
  \texttt{CASE\_063} & Diễn viên An Chinh đã đóng những vai nào? & An Chinh đóng 4 vai: Thị Màu (Quan Âm Thị Kính), Hỷ đồng (Kim Nham), Thiệt Thê (Chu Mãi Thần), Trinh Nguyên (vở Trinh Nguyên) \\
  \hline
  \texttt{CASE\_064} & Diễn viên Tuấn Kha đã đóng những vai nào? & Tuấn Kha đóng 3 vai: Mãng Ông (Quan Âm Thị Kính), Khoèo (Kim Nham), Phù thủy (Kim Nham) \\
  \hline
  \texttt{CASE\_065} & Diễn viên Bá Dũng đã đóng những vai nào? & Bá Dũng đóng 2 vai: Sùng Ông (Quan Âm Thị Kính), Trần Phương (Kim Nham) \\
  \hline
  \texttt{CASE\_066} & Diễn viên Ngọc Minh đã đóng những vai nào? & Ngọc Minh đóng 3 vai: Sùng Ông (Quan Âm Thị Kính), Hề trong trích đoạn Mụ quán - Trần Phương (Kim Nham), Tuần Ty (Chu Mãi Thần) \\
  \hline
  \texttt{CASE\_067} & Diễn viên Thanh Ngoan đã đóng những vai nào? & Thanh Ngoan đóng 2 vai: Sùng Bà (Quan Âm Thị Kính), Đào Huế (Chu Mãi Thần) \\
  \hline
  \texttt{CASE\_068} & Liệt kê toàn bộ vai diễn của diễn viên Tuấn Kha? & Mãng Ông (Quan Âm Thị Kính), Khoèo (Kim Nham), Phù thủy (Kim Nham) — 3 vai ở 2 vở \\
  \hline
  \texttt{CASE\_069} & Tổng kết sự nghiệp diễn xuất của diễn viên An Chinh? & An Chinh đóng 4 vai ở 4 vở: Thị Màu (Quan Âm Thị Kính), Hỷ đồng (Kim Nham), Thiệt Thê (Chu Mãi Thần), Trinh Nguyên (vở Trinh Nguyên) \\
  \hline
  \texttt{CASE\_070} & Tổng kết sự nghiệp của diễn viên Bá Dũng? & Bá Dũng đóng 2 vai: Sùng Ông (QATK), Trần Phương (KN) \\
  \hline
  \texttt{CASE\_071} & Diễn viên Trần Hải đã tham gia những vở nào và vai gì? & Trần Hải đóng vai Phù thủy trong vở Kim Nham (trích đoạn Phù thủy sợ ma) \\
  \hline
  \texttt{CASE\_072} & Diễn viên Ngọc Minh tham gia bao nhiêu vở? Vai gì? & 3 vở: Sùng Ông (QATK), Hề (KN), Tuần Ty (CMT) \\
  \hline
  \texttt{CASE\_073} & Thanh Ngoan tham gia bao nhiêu vở? Liệt kê vai diễn? & 2 vở: Sùng Bà (QATK), Đào Huế (CMT) \\
  \hline
  \texttt{CASE\_074} & Phú Kiên đã tham gia bao nhiêu vở và đóng những vai nào? & Phú Kiên đóng 4 vai ở 4 vở: Phú Ông (Quan Âm Thị Kính), Tuần Ty (Chu Mãi Thần), Từ Thức (vở Từ Thức), Thầy Đồ (vở Trinh Nguyên) \\
  \hline
  \texttt{CASE\_075} & Huy Toàn đã tham gia bao nhiêu vở và đóng những vai nào? & 2 vở: Thiện Sỹ (Quan Âm Thị Kính), Trần Phương (Kim Nham) \\
  \hline
  \texttt{CASE\_076} & Diễn viên nào xuất hiện trong nhiều vở nhất? & An Chinh và Phú Kiên là hai diễn viên tham gia nhiều vở nhất, mỗi người đóng 4 vai ở 4 vở khác nhau \\
  \hline
  \texttt{CASE\_077} & Liệt kê những diễn viên đã tham gia từ 2 vở chèo trở lên? & An Chinh (4 vở), Phú Kiên (4 vở), Ngọc Minh (3 vở), Tuấn Kha (2 vở), Bá Dũng (2 vở), Huy Toàn (2 vở), Thanh Ngoan (2 vở) \\
  \hline
  \texttt{CASE\_078} & Hãy liệt kê những vở chèo cổ tiêu biểu? & 7 vở: Quan Âm Thị Kính, Kim Nham, Trương Viên, Lưu Bình - Dương Lễ, Chu Mãi Thần, Từ Thức, Trinh Nguyên \\
  \hline
  \texttt{CASE\_079} & Vở chèo nào có nhiều nhân vật nhất? & Quan Âm Thị Kính có nhiều nhân vật nhất (Thị Kính, Thị Màu, Thiện Sỹ, Sùng Bà, Sùng Ông, Nô, Mãng Ông, Phú Ông, Lý Trưởng, Mẹ Mõ) \\
  \hline
  \texttt{CASE\_080} & Vở chèo nào có nhiều trích đoạn nhất? & Quan Âm Thị Kính có nhiều trích đoạn nhất: 5 trích đoạn (Cắt râu, Lý trưởng - Mẹ Mõ, Mầu - Nô - Phú Ông, Thị Mầu lên chùa, Vu quy) \\
  \hline
  \texttt{CASE\_081} & Vở chèo nào có nhiều diễn viên tham gia nhất? & Quan Âm Thị Kính với hơn 13 diễn viên \\
  \hline
  \texttt{CASE\_082} & So sánh chủ đề của vở Quan Âm Thị Kính và Kim Nham? & QATK: oan khuất, lòng từ bi Phật giáo. Kim Nham: tình yêu trắc trở, giả điên vì tình. \\
  \hline
  \texttt{CASE\_083} & So sánh nhân vật nữ chính của các vở chèo cổ? & Thị Kính (nhẫn nhịn, tu hành), Súy Vân (phản kháng, giả điên), Châu Long (hy sinh, thủy chung), Thị Phương (hiếu thảo, chịu đựng), Đào Huế (phụ bạc chồng nghèo) \\
  \hline
  \texttt{CASE\_084} & Những vở chèo cổ nào có chủ đề oan khuất? & Có 2 vở: Quan Âm Thị Kính (Thị Kính bị oan), Trương Viên (Thị Phương bị nghi ngờ) \\
  \hline
  \texttt{CASE\_085} & Những vở chèo nào có yếu tố hài (nhân vật Hề)? & Quan Âm Thị Kính, Kim Nham, Trương Viên đều có nhân vật Hề \\
  \hline
  \texttt{CASE\_086} & Những vở chèo nào có chủ đề tình yêu bị cản trở? & Kim Nham (Súy Vân-Kim Nham), Lưu Bình Dương Lễ (Lưu Bình-Châu Long), Từ Thức (Từ Thức-Giáng Hương) \\
  \hline
  \texttt{CASE\_087} & Nhân vật Hề xuất hiện trong những vở chèo nào? & Nhân vật Hề xuất hiện trong Kim Nham (trích đoạn Mụ quán - Trần Phương), Trương Viên (Lớp Tiên Nữ - Đoàn tụ) và Từ Thức (Hề theo Thầy) \\
  \hline
  \texttt{CASE\_088} & Tổng cộng có bao nhiêu trích đoạn trong các vở chèo cổ và phân bổ thế nào? & Quan Âm Thị Kính: 5 trích đoạn; Kim Nham: 4 trích đoạn; Lưu Bình - Dương Lễ: 2 trích đoạn; Chu Mãi Thần, Trương Viên, Từ Thức, Trinh Nguyên: mỗi vở 1 trích đoạn. Tổng cộng 15 trích đoạn. \\
  \hline
  \texttt{CASE\_089} & Có bao nhiêu diễn viên và bao nhiêu nhân vật tổng cộng trong các vở chèo cổ tiêu biểu? & Khoảng 41 diễn viên và 31 nhân vật trải đều trên 7 vở chèo cổ \\
  \hline
  \texttt{CASE\_090} & Nhân vật loại Đào trong chèo là gì? Có những ai là đại diện tiêu biểu? & Đào là vai nữ chính diện: Thị Kính, Châu Long, Thị Phương. Đào lệch: Thị Màu, Súy Vân \\
  \hline
  \texttt{CASE\_091} & Phân tích nhân vật Thị Màu: tính cách, hành động và ý nghĩa? & Thị Màu là nhân vật Đào lệch, táo bạo, phá vỡ khuôn phép phong kiến. Tượng trưng cho khát vọng tự do của người phụ nữ. \\
  \hline
  \texttt{CASE\_092} & Phân tích nhân vật Thị Kính: oan khuất và lòng nhẫn nhịn? & Thị Kính là biểu tượng của đức nhẫn nhịn và lòng từ bi. Bị vu oan 3 lần nhưng vẫn giữ lòng trong sạch, cuối cùng được giải oan. \\
  \hline
  \texttt{CASE\_093} & Phân tích hình tượng người phụ nữ hy sinh trong các vở chèo cổ? & Châu Long (hy sinh 3 năm nuôi bạn chồng), Thị Phương (khoét mắt cho mẹ chồng), Thị Kính (chịu oan suốt đời tu hành) \\
  \hline
  \texttt{CASE\_094} & So sánh số phận bi kịch của Súy Vân và Thị Kính? & Súy Vân bị dụ dỗ giả điên, bị Trần Phương bỏ. Thị Kính bị oan khuất, bị đuổi khỏi nhà, bị đổ oan thai nhi. Cả hai đều là nạn nhân của xã hội. \\
  \hline
  \texttt{CASE\_095} & Triết lý Phật giáo thể hiện như thế nào qua nhân vật Thị Kính? & Thị Kính thể hiện lòng từ bi, nhẫn nhịn, và tin vào nhân quả của Phật giáo. Tên Kính Tâm (kính trọng tâm Phật) thể hiện rõ điều này. \\
  \hline
  \texttt{CASE\_096} & Vai trò của nhân vật phản diện trong các vở chèo cổ? & Sùng Bà (vu oan Thị Kính), Trần Phương (dụ dỗ Súy Vân), Đào Huế (phụ bạc Chu Mãi Thần) — đại diện cho bất công xã hội \\
  \hline
  \texttt{CASE\_097} & Liệt kê các mối quan hệ vợ-chồng trong các vở chèo cổ và kết cục của họ? & Thị Kính-Thiện Sỹ (ly hôn vì vu oan), Trương Viên-Thị Phương (đoàn tụ), Châu Long-Dương Lễ (đoàn tụ), Súy Vân-Kim Nham (tan vỡ), Đào Huế-Chu Mãi Thần (ly hôn) \\
  \hline
  \texttt{CASE\_098} & Nghệ thuật chèo phân loại nhân vật theo những loại gì? & Chèo có 4 loại chính: Đào (nữ chính), Kép (nam chính), Hề (hài), Lão/Cụ (nhân vật lớn tuổi) \\
  \hline
  \texttt{CASE\_099} & Tổng quan toàn bộ nhân vật và diễn viên trong vở Quan Âm Thị Kính? & Vở Quan Âm Thị Kính có 10 nhân vật và hơn 13 diễn viên, 5 trích đoạn. Thị Kính (Thảo Hiền, Hồng Thắm), Thị Màu (Vân Quyền, An Chinh), Thiện Sỹ (Huy Toàn, Trần Xuân Tài), ... \\
  \hline
  \texttt{CASE\_100} & Cho biết tổng quan về các vở chèo cổ tiêu biểu: số vở, nhân vật, diễn viên, trích đoạn? & Gồm 7 vở chèo cổ tiêu biểu, 31 nhân vật, 41 diễn viên, 15 trích đoạn. Vở lớn nhất là Quan Âm Thị Kính. \\
  \hline
\end{longtable}
